\section{Incidents: wrong claims and what caught them}
\label{sec:incidents}

We recorded every case during the final phase in which a statement that had been written or was
about to be written turned out to be wrong or misleading. Table~\ref{tab:incidents} lists them with
the mechanism that detected each. We group the mechanisms into three kinds: \emph{artefact
checks} (a validator or an experiment built to test the statement), \emph{agent cross-checks} (the
agent noticing an inconsistency between two things it had itself produced), and \emph{human
judgment} (the researcher looking at the output and objecting).

\begin{table}[!htbp]
\centering
\caption{Incidents in the final phase. ``Detected by'' is the mechanism that first raised the
problem; A = artefact check (validator or purpose-built experiment), C = agent cross-check, H =
human judgment.}
\label{tab:incidents}
\scriptsize
\begin{tabularx}{\linewidth}{@{}P{0.02\linewidth}P{0.27\linewidth}P{0.25\linewidth}P{0.15\linewidth}Y@{}}
\toprule
\# & Wrong or misleading statement & Cause & Detected by & Resolution \\
\midrule
1 & The manuscript was sound but unreadable: hundreds of ``diagnostic'', ``non-claim'' and
``boundary'' sentences. & Every validator-pinned phrase had been written to satisfy the pin. & H
(``the paper quality feels poor'') & Full rewrite; ledger frozen; one number-checking validator. \\
2 & Orientation error converged at order 2.2 while position, velocity and angular velocity
converged at order 6.2. & Geodesic angle computed as $2\arccos(p\cdot q)$, whose resolution floor
is $3\times10^{-8}$ rad. & C (one of four quantities disagreed) & Chord formula
$2\arcsin(\lVert p\mp q\rVert/2)$; orientation order 6.28. \\
3 & Reference runs on the chain failed at $h=0.1/256$ over $T=1$; the plan said $T=1$ and
$T=20$. & The benchmark's frozen sliding direction degenerates at $t\approx0.605$ ($n_1\cdot a_1\to0$),
a modelling limit, not a solver failure. & A (Newton residual $6\times10^{5}$) then H (chose to
keep the model and stop at $T=0.5$) & All chain experiments on $[0,0.5]$; limit stated in the
paper. \\
4 & Driven single pendulum: velocity error rose and fell erratically between $10^{-15}$ and
$4\times10^{-9}$ while position showed order 6. & The harness's Newton predictor is the analytic
solution and the stage Jacobian has condition number $10^{10}$--$10^{13}$; velocities carry
amplified roundoff. & H (``why does the top-left panel look strange?'') & Velocity not reported for
that mechanism; the panel labelled as a reconstruction check; mechanism explained. \\
5 & The four ASME mechanisms were planned as four convergence tests. & Three of the four are
kinematically driven; their motion is fixed by the drive. & A (errors at $10^{-14}$ for all $h$) &
Only the double pendulum is a dynamics test; the others are consistency rows. \\
6 & The public-code comparison at $T=3$ showed the public methods with $O(1)$ errors; suspicion of
a model mismatch. & The double pendulum amplifies perturbations over $T=3$. & C (added a $T=0.1$
alignment run: first-order decrease, same model) & Comparison kept with the alignment check
reported. \\
7 & ``The $h^7$-scaled stopping rule with $c_\eta=10^{4}$ covers the reported grids.'' & With the
fixed tolerance $10^{-11}$, the accepted residual over $h^7$ reaches $10^{6}$ at the finest
step. & A (purpose-built experiment E7) & Statement replaced by the measured ratios and a
discussion of the pessimistic bound. \\
8 & The single pendulum ``converges at order 6 with a floor'' was fitted over roundoff points,
giving 3.6. & Fit used a zero floor for the analytic reference. & C & Roundoff floor $10^{-15}$;
fit 6.0 over the resolved points. \\
9 & ``All results generated from a single implementation path.'' & Four harnesses assemble the
same stage system for four mechanisms. & C (self-review) & Reworded; assemblies listed. \\
10 & After the rewrite the frozen ledger reported stale line counts and lost markers. & Its
snapshot records pinned the sizes of scripts that had changed. & A (frozen validators) & Builder
sequence rerun; both frozen chains green. \\
\bottomrule
\end{tabularx}
\end{table}

\FloatBarrier
\paragraph{What the artefact checks caught}
Incidents 3, 5, 7 and 10 were caught by things that fail loudly: a Newton solve that does not
converge, a table full of $10^{-14}$, a purpose-built measurement, a validator. Incident~7 is the
instructive one. The analysis section stated that the fixed-tolerance runs satisfy the solver
hypothesis with a modest constant; the agent wrote an experiment (E7) whose only purpose was to
measure that constant, and the measurement contradicted the text by six orders of magnitude on the
finest grid. The text was rewritten to report the measurement. Where experiments are cheap and the
norm is to measure a claimed constant rather than argue for it, a class of overclaims becomes data.

\paragraph{What the agent's cross-checks caught}
Incidents 2, 6, 8 and 9 were caught because the agent had produced two things that had to agree
and did not: four error measures of which one behaved differently, a comparison whose magnitude
called for a sanity check at a shorter horizon, a fit whose value contradicted the visible slope,
and a sentence in the introduction that the methods section did not support. These are the same
habit that produced the central theorem (Section~\ref{sec:trajectory}): redundancy that the agent is
required to reconcile.

\paragraph{What only the human caught}
Incidents 1 and 4, and the decision in Incident 3, needed a person. Incident 1 is a judgment of
quality that no validator expressed and the agent did not have: the manuscript satisfied every check
and was bad. Incident 3 was a modelling decision (keep the benchmark as implemented and stop at
$T=0.5$, or change the joint model) that the agent escalated with three options and a
recommendation to change the model; the human chose to keep it, and the limit was documented.
Incident 4 is the one we find most important. The agent had seen the non-monotone velocity curve,
written a plausible and wrong explanation (a tolerance artefact of the harness), and moved on. The
human's objection was not technical (``it looks strange''). The investigation it triggered found
that the harness starts Newton at the exact solution and that the Jacobian is ill-conditioned by ten
to thirteen orders, a fact about the benchmark that no validator could have known to check. Looking
at figures is not a formality.
